\documentclass[11pt]{article}
\usepackage{preamble}

\newcommand{\IconCheck}{[CHECK]}
\newcommand{\IconMic}{[MIC]}
\newcommand{\IconClipboard}{[CLIPBOARD]}

\newcommand{\phasetag}[1]{%
  {\small\itshape Tag: #1\par}
}

\newcommand{\phaseitems}[1]{%
  \begin{itemize}[leftmargin=1.5em,itemsep=2pt,topsep=2pt]
  #1
  \end{itemize}%
}

\newcommand{\teacheranswerbox}[2]{%
  \begin{callout}{#1}{calloutgreen}
  #2
  \end{callout}
}

\begin{document}

\packetheading{Lesson 3-5 \textperiodcentered\ Period 2 \textperiodcentered\ Teacher Edition}{Multiplicity and Factored-Form Graphing}

\sectionbanner{OBJECTIVES}
{\small
\renewcommand{\arraystretch}{1.25}
\noindent\begin{tabular}{|p{1.8in}|p{4.8in}|}
\hline
\textbf{Math Objective} & Use the multiplicity of each zero to determine whether the graph of a polynomial crosses the x-axis or is tangent to (turns back from) it at each zero. \\ \hline
\textbf{Language Objective} & Students will use the frame ``The zero \(x=\) \_\_\_ has multiplicity \_\_\_, so the graph \_\_\_ the x-axis.'' to describe each zero's graph behavior. \\ \hline
\textbf{Essential Understanding} & A factored-form polynomial reveals its zeros directly. The multiplicity of each zero --- odd or even --- controls whether the graph crosses or is tangent to the x-axis at that zero. \\ \hline
\end{tabular}
}

\sectionbanner{TOPIC GOALS / ESSENTIAL QUESTION / MATERIALS}
{\small
\renewcommand{\arraystretch}{1.25}
\noindent\begin{tabular}{|p{1.8in}|p{4.8in}|}
\hline
\textbf{Topic Goals} & Students apply the multiplicity rule to four factored-form polynomials, one per practice item. Every item is drawn from Savvas Lesson 3-5. \\ \hline
\textbf{Essential Question} & How does the multiplicity of each zero determine the graph's behavior at that x-intercept? \\ \hline
\textbf{Materials} & Pencils; student packet; calculator (optional). No Chromebooks required for Period 2. \\ \hline
\end{tabular}
}

\begin{callout}{\IconPin\ PRE-FLIGHT NOTE (read before class)}{calloutpink}
Explore is 5 bank items: Try It 2a \(\rightarrow\) Try It 2b \(\rightarrow\) \#14 \(\rightarrow\) \#15 \(\rightarrow\) \#16. Paced for \(\sim\)6--7 min per item.\\
If pairs finish early, push them to the Optional Reinforcement: Practice \#11, then Practice \#12.\\
If pace slows, cut \#16 first --- the core multiplicity arc (Try It 2a \(\rightarrow\) Try It 2b \(\rightarrow\) \#14 \(\rightarrow\) \#15) is the minimum viable set.
\end{callout}

\phasetag{bridge from Monday}
\frameworkphaseheader{Do Now}{1}{5}{%
\phaseitems{
  \item Hand out Do Now at door.
  \item Circulate during 3 min work.
  \item Cold-call 1 student to share one zero at minute 4.
}%
}{%
\phaseitems{
  \item Read each factor; write zeros directly.
  \item Use sentence frame on the sheet.
}%
}{%
\phaseitems{
  \item What does it mean for \((x-a)\) to be a factor?
  \item If \((x+6)\) is a factor, what zero does that give you?
}%
}{Solo teacher: stand at door, hand out Do Now as students enter. Circulate for \(\sim\)3 min. Cold-call 1 student to share one zero.}

\bankitem{Do Now (3-5-savvas-q13)}{%
Sketch the graph of the function by finding the zeros: \(g(x)=(x+3)(x-1)(x-4)\).
}{}

\teacheranswerbox{\IconCheck\ ANSWER KEY}{%
Zeros: \(x=-3\), \(x=1\), \(x=4\). Each factor \((x+3)\), \((x-1)\), \((x-4)\) equals zero when \(x\) takes the opposite sign of its constant term.
}

\begin{callout}{\IconTarget\ SENTENCE FRAME (for your answer)}{calloutyellow}
``The zeros are \(x=\) \_\_\_, \(x=\) \_\_\_, and \(x=\) \_\_\_ because each factor equals zero.''
\end{callout}

\sectionbanner{LAUNCH --- Multiplicity Rule (teacher models on screen)}

\phasetag{teacher models Example 2}
\frameworkphaseheader{Launch}{2}{12}{%
\phaseitems{
  \item Project Savvas Example 2. Think aloud: count each factor; odd \(\rightarrow\) crosses, even \(\rightarrow\) tangent.
  \item Pause 2\(\times\) for 30-sec partner-checks.
  \item Do NOT model Try-Its --- release to Explore.
}%
}{%
\phaseitems{
  \item Watch the modeled example silently first.
  \item Turn-and-talk during each 30-sec pause.
  \item Copy the Multiplicity Rules card into their page margin if desired.
}%
}{%
\phaseitems{
  \item Say: ``Count the factor. Odd or even?'' for each zero.
  \item What's the difference between ``crosses'' and ``tangent to''?
  \item Why does the graph turn back at even multiplicity? (Sign doesn't flip.)
}%
}{Project Savvas Example 2. Work it aloud with think-aloud voice. Pause 2\(\times\) for 30-sec partner-check: ``Tell your neighbor what multiplicity this zero has.''}

\begin{callout}{\IconBook\ MULTIPLICITY RULES (keep printed on your page)}{calloutgreen}
\textbullet\ If a factor appears ONCE or an ODD number of times, the zero has ODD multiplicity \(\rightarrow\) the graph CROSSES the x-axis at that zero.\\
\textbullet\ If a factor appears an EVEN number of times, the zero has EVEN multiplicity \(\rightarrow\) the graph is TANGENT to the x-axis (touches and turns back) at that zero.\\
\textbullet\ Example: \(f(x)=(x-2)^3(x+1)^2\) \(\rightarrow\) \(x=2\) has multiplicity 3 (crosses); \(x=-1\) has multiplicity 2 (tangent).
\end{callout}

\begin{callout}{\IconMic\ TEACHER SCRIPT}{calloutblue}
1. ``Watch me do this first one. I'll read each factor, count how many times it appears, and decide if the graph crosses or is tangent.''\\
2. For \(f(x)=(x-2)^3(x+1)^2\): ``At \(x=2\), the factor appears 3 times --- odd --- so the graph CROSSES. At \(x=-1\), the factor appears 2 times --- even --- so the graph is TANGENT.''\\
3. ``What's one thing I just did that I want you to copy?'' --- cold-call.\\
4. Do NOT do Try-Its for them. Set them up, release to TPS.
\end{callout}

\sectionbanner{EXPLORE --- Think-Pair-Share (35+ min)}
\textit{Work with a partner. Use the multiplicity rules above for every problem.}

\phasetag{Think-Pair-Share + Practice}
\frameworkphaseheader{Explore}{2}{33}{%
\phaseitems{
  \item 4 laps of circulation (\(\sim\)8 min each).
  \item Lap 1: notice stalls. Lap 2: ask purposeful Qs. Lap 3: check answers. Lap 4: invite 1 pair to share.
  \item Do not give answers. Press pairs to justify with the rule card.
}%
}{%
\phaseitems{
  \item Work in pairs through Try-It 2a, 2b, then Practice \#14, \#15, \#16.
  \item For each zero, write multiplicity + crosses/tangent before sketching.
  \item Justify with sentence frame aloud to partner before writing.
}%
}{%
\phaseitems{
  \item What's the multiplicity of each zero?
  \item Does the graph cross or is it tangent at each zero?
  \item How does the end behavior match the leading coefficient + degree?
}%
}{Solo teacher: 4 laps of circulation (\(\sim\)8 min each). Lap 1: notice who's stalling. Lap 2: ask purposeful Qs. Lap 3: check answers. Lap 4: invite 1 pair to share at board.}

\textbf{Bank-item labels (in order):}
\begin{enumerate}[leftmargin=1.6em,itemsep=2pt,topsep=2pt]
  \item Try It 2a
  \item Try It 2b
  \item Practice \#14
  \item Practice \#15
  \item Practice \#16
\end{enumerate}

\bankitem{Try It 2a}{%
Describe the behavior of the graph of the function at each of its zeros: \(f(x)=x(x+4)(x-1)^4\).
}{}

\teacheranswerbox{\IconCheck\ ANSWER / ANTICIPATED TALK}{%
\textbf{Answer key:} Graph crosses the x-axis at zeros \(0\) and \(-4\) (odd multiplicity 1); turning point at \(x=1\) (even multiplicity 4).
}

\bankitem{Try It 2b}{%
Describe the behavior of the graph of the function at each of its zeros: \(f(x)=(x^2+9)(x-1)^5(x+2)^2\).
}{}

\teacheranswerbox{\IconCheck\ ANSWER / ANTICIPATED TALK}{%
\textbf{Answer key:} Real zeros only: graph crosses at \(x=1\) (odd multiplicity 5); turning point at \(x=-2\) (even multiplicity 2). Factor \(x^2+9\) has no real zeros.
}

\bankitem{Practice \#14}{%
Find the zeros of the function, and describe the behavior of the graph at each zero: \(f(x)=x^3-8x^2+16x\).
}{}

\teacheranswerbox{\IconCheck\ ANSWER / ANTICIPATED TALK}{%
\textbf{Answer key:} Factor: \(x(x-4)^2\). Zeros: \(0, 4\). The graph crosses the x-axis at \(0\) (multiplicity 1) and is tangent (touches and turns back) at \(4\) (multiplicity 2).
}

\bankitem{Practice \#15}{%
Find the zeros of the function, and describe the behavior of the graph at each zero: \(g(x)=x^3-x^2-25x+25\).
}{}

\teacheranswerbox{\IconCheck\ ANSWER / ANTICIPATED TALK}{%
\textbf{Answer key:} Factor by grouping: \(x^2(x-1)-25(x-1)=(x-1)(x^2-25)=(x-1)(x-5)(x+5)\). Zeros: \(-5, 1, 5\). The graph crosses the x-axis at \(-5\), \(1\), and \(5\) (all simple zeros).
}

\bankitem{Practice \#16}{%
Find the zeros of the function, and describe the behavior of the graph at each zero: \(f(x)=9x^4-40x^2+16\).
}{}

\teacheranswerbox{\IconCheck\ ANSWER / ANTICIPATED TALK}{%
\textbf{Answer key:} Biquadratic: let \(u=x^2\) \(\rightarrow\) \(9u^2-40u+16=(9u-4)(u-4)=0\) \(\rightarrow\) \(u=4/9\) or \(u=4\) \(\rightarrow\) \(x=\pm 2/3\) or \(x=\pm 2\). Zeros: \(\pm 2/3, \pm 2\). The graph crosses the x-axis at each zero (all simple).
}

\sectionbanner{SHARE / SUMMARY (5 min)}

\phasetag{academic conversation + bridge prompt}
\frameworkphaseheader{Share / Summary}{2}{5}{%
\phaseitems{
  \item Cold-call 2 pairs to read their sentence-frame sentence.
  \item Write the rule sentence on board (student-generated).
  \item Post the Wednesday bridge prompt visibly.
}%
}{%
\phaseitems{
  \item One student per pair reads their sentence aloud.
  \item Class signals thumbs up/sideways.
  \item Write the bridge prompt's key term (conjugate pair) in packet margin.
}%
}{%
\phaseitems{
  \item What's a sentence that captures today's rule?
  \item For Wednesday: if a cubic has zeros 2, \(1+\sqrt{3}\), and \(1-\sqrt{3}\) --- what does that say about its coefficients?
}%
}{Cold-call 2 pairs. Write one sentence on board. Leave the bridge prompt posted for Wednesday Do Now.}

\begin{sentenceframebox}
\textbf{Sentence frames}

Pair share: ``For \(f(x)=\) \_\_\_, the zero \(x=\) \_\_\_ has multiplicity \_\_\_, so the graph \_\_\_ the x-axis at that zero.''

Whole class: one pair at random reads their sentence. Class signals thumbs up/sideways.
\end{sentenceframebox}

\begin{callout}{\IconSearch\ BRIDGE PROMPT (post for Wednesday Do Now)}{calloutpeach}
``A cubic has zeros 2, \(1+\sqrt{3}\), and \(1-\sqrt{3}\). What is the nature of its coefficients --- rational, irrational, or complex? Why?'' (Directly prepares Topic 3 Assessment Q\#2.)
\end{callout}

\sectionbanner{EXIT TICKET}

\phasetag{summary recap (not graded)}
\frameworkphaseheader{Exit Ticket}{1-2}{2}{%
\phaseitems{
  \item Collect at door as students exit.
  \item Scan for confusion flags; plan tomorrow's Do Now emphasis accordingly.
}%
}{%
\phaseitems{
  \item Complete 3 sentence stems.
  \item Be honest about confusion --- this is not graded.
}%
}{%
\phaseitems{
  \item Did students write a complete sentence for \#1?
  \item Did anyone flag confusion in \#3?
}%
}{Collect at door. Scan flagged items tonight. No formal grade.}

\begin{callout}{\IconClipboard\ EXIT TICKET --- Student Form}{calloutyellow}
Today I learned that \_\_\_.\\
The rule I want to remember is \_\_\_.\\
One thing I'm still unsure about is \_\_\_.
\end{callout}

\sectionbanner{OPTIONAL REINFORCEMENT (not collected)}
\textit{If you finish early or want extra practice before Wednesday. Answers on Desmos or ask for a check.}

\bankitem{Optional \#1  (Savvas Practice)}{%
At what points do the graphs of \(f(x)=x^3-2x^2-16x+20\) and \(g(x)=-12\) intersect?
}{}

\teacheranswerbox{\IconCheck\ ANSWER / ANTICIPATED TALK}{%
\textbf{Answer key:} Solve \(x^3-2x^2-16x+20=-12\) \(\rightarrow\) \(x^3-2x^2-16x+32=0\) \(\rightarrow\) \((x-2)(x+4)(x-4)=0\). Intersection points: \((-4,-12)\), \((2,-12)\), \((4,-12)\).
}

\bankitem{Optional \#2  (Savvas Practice)}{%
Sketch the graph of the function by finding the zeros: \(f(x)=3x^3-9x^2-12x\).
}{}

\teacheranswerbox{\IconCheck\ ANSWER / ANTICIPATED TALK}{%
\textbf{Answer key:} Factor: \(3x(x+1)(x-4)\). Zeros: \(x=-1, 0, 4\). Cubic with positive leading coefficient, three simple zeros (graph crosses at each).
}

\clearpage

\begin{callout}{\IconPuzzle\ MODIFICATIONS / SUPPORT FOR IEP STUDENTS}{calloutiep}
\textbullet\ Provide the Multiplicity Rule card as a sticker/cut-out on their desk.\\
\textbullet\ Allow color-coding each zero on their factored form before answering.\\
\textbullet\ Accept verbal sentence-frame completions in lieu of written.
\end{callout}

\begin{callout}{\IconGlobe\ SUPPORTS FOR ELL STUDENTS}{calloutell}
\textbullet\ Sentence frames printed on student packet (don't skip).\\
\textbullet\ Word card: ``multiplicity'' = how many times a factor appears.\\
\textbullet\ ``crosses'' = goes through; ``tangent'' = touches and turns back (Savvas terminology).\\
\textbullet\ Pair ELL students with English-dominant partners for TPS whenever possible.
\end{callout}

\begin{callout}{\IconClock\ PERIOD F EXTENSION (65-min class)}{calloutpink}
Period F has an extra 10 min. Use Explore +10 min (not Launch +10).\\
Suggested add: after Practice \#16, pair students do Practice \#11 or \#12 from the Optional Reinforcement page.\\
Do NOT add a second Launch example --- keeping the rhythm simple.
\end{callout}

\end{document}
